\documentclass[11pt]{article}
\usepackage[a4paper,margin=1in]{geometry}
\usepackage{fontspec}
\usepackage[boldfont=false]{xeCJK}
\setCJKmainfont{FandolSong-Regular.otf}
\usepackage{amsmath}
\usepackage{amssymb}
\usepackage{array}
\usepackage{booktabs}
\usepackage{graphicx}
\usepackage{adjustbox}
\usepackage{enumitem}
\setlist[itemize]{topsep=2pt,partopsep=0pt,parsep=0pt,itemsep=2pt,leftmargin=1.4em}
\setlist[enumerate]{topsep=2pt,partopsep=0pt,parsep=0pt,itemsep=2pt,leftmargin=1.6em}
\usepackage{parskip}
\usepackage{fvextra}
\fvset{breaklines=true,breakanywhere=true,fontsize=\small}
\setlength{\parindent}{0pt}

\begin{document}

\begin{center}
  {\LARGE\bfseries Business and its environment (A Level)}\\[0.35em]
  {\large A Level Business}
\end{center}
\vspace{0.6em}

\section*{External influences on business}

A business cannot control everything. \textbf{External influences} 外部影响 are forces outside the firm that affect how it works. The firm must watch them and adapt. We group them into economic factors and wider (political, legal, social and other) factors.

\section*{The business cycle}

The \textbf{business cycle} 经济周期 is the way the whole economy rises and falls over time. It has four stages:

\begin{itemize}
\item \textbf{boom} 繁荣 — fast growth; high demand, but rising costs and prices.

\item \textbf{recession} 衰退 — falling demand and output; sales drop and some firms close.

\item \textbf{slump} 萧条 — a deep, long recession with very low demand.

\item \textbf{recovery} 复苏 — demand and output start to rise again.

\end{itemize}

In a boom, firms expand; in a recession, they cut costs, hold less stock, and may aim only to survive.

\section*{Interest rates}

The \textbf{interest rate} 利率 is the price of borrowing money, set as a percentage. When interest rates \textbf{rise}:

\begin{itemize}
\item loans and mortgages cost more, so customers spend less.

\item firms with loans pay more, which lowers their profit.

\item new investment is less likely, because borrowing is dear.

\end{itemize}

When interest rates fall, the opposite happens and spending usually rises.

\section*{Exchange rates}

The \textbf{exchange rate} 汇率 is the price of one currency in terms of another. It changes all the time.

\begin{itemize}
\item if the local currency rises in value (\textbf{appreciation} 升值), exports become dearer abroad and imports become cheaper.

\item if the local currency falls in value (\textbf{depreciation} 贬值), exports become cheaper abroad and imports become dearer.

\end{itemize}

So an \textbf{exporter} 出口商 often gains from a weaker currency, while an \textbf{importer} 进口商 often gains from a stronger one.

\section*{Inflation}

\textbf{Inflation} 通货膨胀 is a general rise in prices over time. High inflation raises a firm's costs (materials and wages), makes planning harder, and can cut what customers can afford to buy. Low, steady inflation is usually best for business.

\section*{Taxation and unemployment}

\textbf{Taxation} 税收 is money the government takes from people and firms.

\begin{itemize}
\item a \textbf{direct tax} 直接税 is taken from income or profit (e.g. income tax, corporation tax).

\item an \textbf{indirect tax} 间接税 is added to spending (e.g. sales tax on goods).

\end{itemize}

Higher taxes leave people and firms with less money to spend, so demand falls.

\textbf{Unemployment} 失业 is the number of people who want work but cannot find it. High unemployment lowers demand, but it also makes workers easier and cheaper to hire.

\section*{Political, legal and social influences}

\begin{itemize}
\item \textbf{political} 政治 — government decisions, trade rules and stability affect how firms plan.

\item \textbf{legal} 法律 — firms must obey \textbf{legislation} 法规 on employment, consumer safety, competition and the environment. Breaking the law brings fines and bad publicity.

\item \textbf{social} 社会 — changes in society, such as ageing populations or new tastes, change what people buy.

\end{itemize}

\section*{Technological, environmental and ethical influences}

\begin{itemize}
\item \textbf{technological} 技术 — new technology can create new products and cut costs, but firms must keep up or fall behind.

\item \textbf{environmental} 环境 — firms face pressure and rules to cut pollution and waste, and to act in a sustainable way.

\item \textbf{ethical} 道德 — acting in a morally right way (fair pay, honest selling) builds trust, even if it costs more.

\end{itemize}

Together these wider factors are often studied with a \textbf{PESTLE analysis} 宏观环境分析, which checks Political, Economic, Social, Technological, Legal and Environmental forces.

\section*{What strategy means}

A \textbf{strategy} 战略 is a long-term plan to reach the firm's main objectives. \textbf{Strategic management} 战略管理 is the work of setting that plan, putting it into action, and checking it. Good strategy aims to build a \textbf{competitive advantage} 竞争优势 — a reason customers choose you over rivals, such as lower cost or a better product.

\section*{SWOT analysis}

A \textbf{SWOT analysis} SWOT分析 studies a firm's position under four headings:

\begin{center}
\adjustbox{max width=\linewidth}{%
\begin{tabular}{ll}
\toprule
\textbf{Internal (inside the firm)} & \textbf{External (outside the firm)} \\
\midrule
\textbf{strengths} 优势 — what it does well & \textbf{opportunities} 机会 — chances to grow \\
\textbf{weaknesses} 劣势 — what it does poorly & \textbf{threats} 威胁 — dangers it faces \\
\bottomrule
\end{tabular}}
\end{center}

A firm should use its strengths to take opportunities, and fix or protect against its weaknesses and threats.

\section*{Other planning tools}

Firms use other tools to make decisions, such as decision trees (weighing choices by their likely results) and the matrices you meet later in marketing and operations. Every planning tool turns information about the business and its environment into a clearer choice.


\end{document}
